% =====================================================================
\section{Introducci\'on}\label{sec:introduccion}
% =====================================================================
% Que debe ir aqui (aprox. 1 pagina):
%  - Problema de la organizacion en 1 parrafo (con datos de lo observado).
%  - Proposito de este documento y su lugar en el PFC.
%  - Marco de referencia: por que se identifican partes interesadas y el
%    limite del sistema ANTES de definir requisitos \cite{iso29148}; como
%    se eligen tecnicas de elicitacion segun la fuente y el tipo de
%    conocimiento \cite{pohl2025,ireb2026}.
%  - Organizacion del documento (una frase por seccion, con \ref).
% Redacte en tercera persona, sin relleno.

\pendiente{Redactar la introducci\'on con base en el problema real de la organizaci\'on.}
\pendiente{Marco: partes interesadas y l\'imite del sistema \cite{iso29148}; selecci\'on de t\'ecnicas \cite{pohl2025,ireb2026}.}

El documento se organiza as\'i: la Secci\'on~\ref{sec:organizacion} presenta la organizaci\'on;
la Secci\'on~\ref{sec:vision} define la visi\'on, los objetivos y el alcance;
la Secci\'on~\ref{sec:partes} identifica a las partes interesadas;
la Secci\'on~\ref{sec:contexto} describe el contexto del sistema;
la Secci\'on~\ref{sec:elicitacion} detalla el plan de elicitaci\'on;
la Secci\'on~\ref{sec:etica} trata la \'etica y el manejo de datos;
la Secci\'on~\ref{sec:gestion} presenta la gesti\'on del proyecto;
y la Secci\'on~\ref{sec:conclusiones} resume las conclusiones.
